\documentclass{article}

% The official style file. Download
%   https://media.neurips.cc/Conferences/NeurIPS2025/Styles.zip
% and place neurips_2025.sty beside this file. A layout-matched stand-in is
% included in this directory so the page count can be checked; overwrite it
% with the official file before submitting.
\usepackage[final]{neurips_2025}

\usepackage[utf8]{inputenc}
\usepackage[T1]{fontenc}
\usepackage{hyperref}
\usepackage{url}
\usepackage{booktabs}
\usepackage{amsfonts}
\usepackage{amsmath}
\usepackage{nicefrac}
\usepackage{microtype}
\usepackage{xcolor}

\title{Faithful Diffusion-Based Image Super-Resolution}

% TODO: replace with the full author list.
\author{%
  Shambhavi \textcolor{red}{[last name]}, University of California San Diego\\
  \texttt{shambhavi12001@gmail.com}
}

\begin{document}
\maketitle

\section{Research problem}

Conditional diffusion models give the best perceptual quality known for
single-image super-resolution \citep{saharia2023sr3}, sampling from a learned
approximation of the posterior $p(x \mid y)$ over high-resolution images
consistent with a low-resolution observation $y$. At $4\times$ that posterior
is broad, so the sampler commits to detail the measurement does not determine:
a plausible texture, a sharpened edge, a character that looks legible on no
evidence. Perceptually this is an improvement; as a measurement it is
fabrication.

Partial forward diffusion \citep{meng2022sdedit} makes the amount controllable
through one scalar $s \in [0,1]$: noise the bicubic upsample of $y$ to timestep
$t_0 = \lfloor sT \rfloor$, then denoise to zero; $s=0$ returns bicubic, $s=1$ is full
generative sampling. For super-resolution this knob
is essentially unmeasured.

We ask: \textbf{(Q1)} for a fixed model and fixed $y$, as $s$ rises, how do
perceptual quality, distortion, measurement fidelity and hallucination each
move, and does the pattern differ for stochastic texture versus regular
man-made structure? \textbf{(Q2)} does an explicit reconstruction
constraint---projecting each sampler iterate back onto the images consistent
with $y$---\emph{translate} the trade-off curve toward faithfulness, or only
move the operating point along it? \citet{blau2018perception} prove a trade-off
must exist, not that a given sampler sits on its frontier; Q2 asks whether this
one does.

\section{Why this problem is important}

Because the trade-off is unavoidable \citep{blau2018perception} there is no
best operating point, only a choice---made implicitly today,
by picking a model and a sampler, with no map of what the alternatives cost.
The cost of an invented detail is also asymmetric: a blurry reconstruction is
honest about its ignorance, a sharp wrong one is not, and in medical, forensic
and scientific imaging the second failure is the damaging one
\citep{tivnan2024hallucination}. Hallucination has meanwhile acquired a precise
definition: \citet{cohen2024looks} identify it with posterior uncertainty,
measurable as disagreement among samples sharing $y$, which is what makes a
systematic sweep possible now.

\section{Existing work and its limitations}

SR3 \citep{saharia2023sr3} established conditional diffusion super-resolution
by concatenating the bicubic upsample of $y$ to the denoiser input; it notes
that distortion metrics worsen, but samples at one fixed strength. SDEdit
\citep{meng2022sdedit} introduced the strength knob as a realism--faithfulness
dial, but for stroke-based editing, with no ground truth and no measurement
operator, so it quantifies neither distortion nor hallucination.
Consistency-enforcing samplers---DDNM \citep{wang2023ddnm}, diffusion posterior
sampling \citep{chung2023dps}, ILVR \citep{choi2021ilvr}---show fidelity to
$y$ can be imposed, but each evaluates a single configuration; none asks how the
constraint reshapes the trade-off. Closest to our question,
\citet{tivnan2024hallucination} note in passing that more modest denoising
strengths reduce hallucination, on one medical modality with a single scalar
index, no sweep and no intervention tested; and \citet{cohen2024looks} supply
the uncertainty-based estimator, validated on BSD100 at $4\times$ bicubic,
exactly our protocol, but to compare models rather than to trace one across
strengths.

The field thus has a theorem, isolated point measurements, a knob borrowed from
editing, and one anecdote; missing is a \emph{measured curve} of faithfulness
and hallucination against perceptual quality as continuous functions of strength
within one trained model, plus a controlled test of whether data consistency
moves it.

\section{Proposed solution}

\textbf{Model and knob.} One conditional diffusion model in the SR3 formulation
\citep{saharia2023sr3}: $\epsilon$-prediction DDPM objective
\citep{ho2020ddpm}, bicubic upsample of $y$ concatenated channel-wise to the
noisy target, trained at $4\times$ on the 800 DIV2K training images
\citep{agustsson2017div2k} with our own degradation operator. Exactly one model
is trained; all baselines are deterministic or pretrained.
Sampling is DDIM \citep{song2021ddim} with $\eta=0$ from $t_0 = \lfloor sT
\rfloor$, step spacing fixed at $T/50$, so smaller $s$ means strictly fewer
network evaluations and compute becomes a reported third axis.

\textbf{Measurement}, per $(s, \text{seed}, \text{image})$: LPIPS
\citep{zhang2018lpips}; $Y$-channel PSNR and SSIM \citep{wang2004ssim} with a
4-pixel border crop, per the EDSR protocol \citep{lim2017edsr}; LR consistency,
the PSNR between the bicubic decimation of the output and $y$; the per-pixel
standard deviation across seeds, image-averaged, as
hallucination-as-posterior-uncertainty \citep{cohen2024looks}; and median
wall-clock. The last two need no ground truth.

\textbf{Design.} $s \in \{0, 0.2, 0.4, 0.6, 0.8, 1.0\}$ crossed with five seeds
on BSD100 \citep{martin2001bsds} and Urban100 \citep{huang2015urban100}, a
contrast pair: invented stochastic texture is perceptually cheap, invented
regular structure is not. Set5 and Set14 only confirm the pipeline
reproduces published numbers; baselines are bicubic ($s=0$) and pretrained EDSR
\citep{lim2017edsr}. We then insert DDNM back-projection \citep{wang2023ddnm},
$\hat{x}_0 \leftarrow \hat{x}_0 - H^\dagger(H\hat{x}_0 - y)$, at every step
and repeat the identical sweep, so the constraint is the only difference between
the curves. Analysis is paired per image: bootstrap intervals,
Wilcoxon signed-rank between adjacent strengths with Holm correction
\citep{holm1979simple}, Spearman for monotonicity. Hypotheses: \textbf{(H1)}
LPIPS improves while PSNR and LR consistency degrade monotonically in $s$;
\textbf{(H2)} cross-seed deviation rises with $s$ and correlates with the LPIPS
gain; \textbf{(H3)}
back-projection translates the curve, giving equal LPIPS at higher LR
consistency.

\section{Novelty and significance}

The novelty is the measurement, not a new sampler. We are not aware of work
that traces the perception--faithfulness trade-off as a continuous function of
generative strength within one diffusion model, with hallucination measured as
posterior uncertainty rather than inferred from falling PSNR, and that then
tests whether a data-consistency constraint translates the curve. We do not claim
H1's phenomenon as new---\citet{tivnan2024hallucination} noted it---only its
systematic measurement, with uncertainty quantified and multiple-comparison
control. The product is an operating-point map: for a faithfulness budget in dB
of LR consistency, which $s$ to run and what it costs in LPIPS and in seconds.
It separates blur from fabrication, two failures PSNR scores identically, and
settles whether back-projection is free or the frontier is tight as
\citet{blau2018perception} predict.

\section{Tentative timeline}

\begin{table}[h]
\centering
\scriptsize
\setlength{\tabcolsep}{5pt}
\begin{tabular}{@{}lll@{}}
\toprule
Week & Dates & Deliverable and gate \\
\midrule
1 & Oct 21--27 & Degradation, six metrics, data. \emph{Done: 21/21 checks, 28.9\,dB} \\
2 & Oct 28--Nov 3 & SR3 denoiser, DDIM sampler; overfit-ten and $s{=}1$ gates \\
3 & Nov 4--10 & Train on DIV2K; checkpoint chosen on validation split \\
4 & Nov 11--17 & Strength knob; main sweep ($6{\times}5{\times}200$), resumable \\
5 & Nov 18--24 & Back-projection variant; repeat sweep; statistics \\
6 & Nov 25--Dec 1 & Figures, write-up, buffer \\
\bottomrule
\end{tabular}
\end{table}

\textbf{Fallback.} If training has not converged by week~3 we substitute a
public pretrained checkpoint; if the sweep runs long we drop to three seeds and
BSD100 only. Either outcome of H3 is reportable, so nothing is all-or-nothing.

\bibliographystyle{plainnat}
\bibliography{refs}

\end{document}
